% New deduction, root, JAC1964-20261121. Whole-proof review recorded.
\subsection*{AR-18. A direct comparison of canonical surface heights}

The adelic formula also turns the missing sign into a concrete comparison
between two effective cycles. This gives a second geometric target for
the slope approach; it is not a proof of the comparison.

\begin{proposition}[AR-18]
Let $A/K$ be an abelian threefold over a number field, with a rigidified
symmetric ample line bundle $L$ and its canonical metric. Let $V$ be an
effective rational one-cycle, put
\[
 z=\operatorname{Ce}(V)=\frac{V-[-1]_*V}{2},
 \qquad Lz=0\quad\text{in }\operatorname{CH}^3(A)_{\mathbf Q},
\]
and let $\mu_\pm(x,y)=x\pm y$. Define the effective rational two-cycles
$S_\pm=(\mu_\pm)_*(V\times V)$, with the usual proper-pushforward
multiplicities and zero contribution from a component whose image has
smaller dimension. Their unnormalized canonical intersection heights are
\[
 H_\pm(V)=\frac1{[K:\mathbf Q]}
       \widehat{\deg}(\overline L^3\cdot S_\pm).
\]
Then
\[
 h_K(z,z)=\frac{H_-(V)-H_+(V)}{12h^0(A,L)}.                 \tag{AR.34}
\]
Every primitive class in $\operatorname{CH}^2_{(1)}(A)_{\mathbf Q}$
admits an effective rational representative in this sense. Consequently
semipositivity on that whole space is equivalent to
$H_-(V)\ge H_+(V)$ for every such $V$. The known nonnegativity
$H_\pm(V)\ge0$ does not itself order these two numbers.
\end{proposition}

\begin{proof}
Write $\overline L_i=p_i^*\overline L$ and
$\overline Q=\mu_+^*\overline L-\overline L_1-\overline L_2$.
The theorem of the cube and the canonical rigidifications give the
metric identities
\[
 \mu_\pm^*\overline L=\overline L_1+\overline L_2\pm\overline Q,
 \qquad (1,[-1])^*\overline Q=-\overline Q.
\]
Inversion in either factor therefore makes intersection against an odd
power of $\overline Q$ odd in the corresponding cycle. In particular
\[
 \widehat{\deg}(\overline Q^3\cdot(V\times V))
 =\widehat{\deg}(\overline Q^3\cdot(z\times z)).             \tag{AR.35}
\]
All identities are identities of arithmetic intersections of the chosen
cycles; descent of an arbitrary such intersection to a bare Chow group
is not being assumed.

We need the vanishing of the mixed term. Let $D=L\cdot V$ and let
$u=\operatorname{alb}_A(D)$, where the origin of $A$ defines the
Albanese sum for zero-cycles of any degree. Since $L$ is symmetric,
\[
 \operatorname{alb}_A(Lz)
 =\tfrac12\bigl(\operatorname{alb}_A(D)
                    -\operatorname{alb}_A([-1]_*D)\bigr)=u.
\]
Thus $u=0$ in $A(K)\otimes\mathbf Q$ under the stated primitive
hypothesis. All line bundles and pairings in this paragraph are taken
in the rational adelic Picard group; equivalently, one can first clear
the rational coefficients and then kill the resulting torsion point.

Here is a metric justification for using this ordinary zero-cycle.
Normalize each curve in $V$ and use additivity of the Deligne pairing.
Over the first factor $A$, the relative Deligne pairing of
$\overline Q$ with $\overline L$ along the second curve has underlying
line bundle $Q|_{A\times\{u\}}$. This is the norm description of the
pairing: a divisor for $L$ gives the tensor product of its Poincare
fibres, and the theorem of the cube identifies that product with the
fibre at the Albanese sum. Moreover, pulling back by $[n]$ on the first
factor multiplies this metrized pairing by $n$, since it does so to
$\overline Q$. The rigidification at zero is trivial. Hence its metric
is exactly the canonical Picard metric. Uniqueness follows directly:
the logarithm of a ratio of two such continuous metrics is bounded on
the compact analytic space and satisfies $f([n]x)=n f(x)$, forcing
$f=0$. For $u=0$ in $A(K)\otimes\mathbf Q$ this is the zero class
of the rational adelic Picard group. Applying the
Deligne-pairing projection formula once more gives
\[
 \widehat{\deg}(\overline L_1\overline L_2\overline Q
                   \cdot(V\times V))=0.                   \tag{AR.36}
\]
Also $\overline L_1^2\overline Q$ and
$\overline L_2^2\overline Q$ have zero intersection against
$V\times V$: pair first over the other curve; the degree of its
Poincare restriction is zero, so the projection formula gives the
zero tensor power. These observations apply to singular curves through
normalization and proper projection, and to rational coefficients by
multilinearity.

Expanding the two cubes now gives
\[
 \begin{split}
 [K:\mathbf Q](H_-(V)-H_+(V))
 &=-6\widehat{\deg}\bigl((\overline L_1+\overline L_2)^2
                           \overline Q\cdot(V\times V)\bigr)\\
 &\phantom{={}}-2\widehat{\deg}\bigl(\overline Q^3
                           \cdot(V\times V)\bigr)\\
 &=-2\widehat{\deg}\bigl(\overline Q^3\cdot(z\times z)\bigr).
 \end{split}
\]
The odd part of a one-cycle on a threefold has Beauville grade one:
Fourier sends $\operatorname{CH}^2_{(s)}$ to
$\operatorname{CH}^{1+s}_{(s)}$; negative grades and grades above two
vanish in these degrees, and inversion selects odd $s$.
AR-10 with $m=2$ therefore identifies the last expression with
$12h^0(A,L)[K:\mathbf Q]h_K(z,z)$, proving (AR.34).

Finally represent any grade-one primitive class $z$ by a finite rational
cycle $Z=Z_+-Z_-$ with $Z_+,Z_-$ effective. Set
$V=Z_++[-1]_*Z_-$. It is effective and
$\operatorname{Ce}(V)=\operatorname{Ce}(Z)=z$ in rational Chow groups.
The already proved identity supplies the claimed equivalence.
Nonnegativity of the canonical heights of effective cycles is Zhang,
\emph{Small points and adelic metrics}, Theorem 2.4(b), author manuscript
pp.~11--12; apply it to each component of $S_\pm$ and add. It gives two
separate inequalities, not their difference.\end{proof}

\paragraph{Exact next gate for this route.}
For the complementary Fermat group the comparison is required for all
effective $V$ whose Ceresa class lies in $N$, not only for the Fermat
curve. Neither canonical-height nonnegativity nor the two nef bundles
$\mu_\pm^*\overline L$ produces the missing ordering. The adelic height
line of AR-10 can equally be represented, at the level of arithmetic
degree, by the difference of these two surface Deligne pairings divided
by $12h^0(A,L)$. A Bost proof along this route must provide a comparison
map and its norm bound between those particular pairings.

The exact background is Gu, arXiv:2609.33223v2, Theorem 1.3,
Proposition 5.3 and \S2.2, and Zhang, Theorem 2.4(b),
\url{https://web.math.princeton.edu/~shouwu/publications/smpoint.pdf}.
The relation (AR.34), the primitive mixed-term vanishing, and the
effective representative construction are the deduction here.
